\documentclass[11pt,a4paper]{article}
\usepackage[margin=1in]{geometry}
\usepackage[T1]{fontenc}
\usepackage{lmodern}
\usepackage{microtype}
\usepackage{amsmath,amssymb,mathtools,bm}
\usepackage{booktabs,tabularx,array}
\usepackage{xcolor}
\usepackage{listings}
\usepackage{tikz}
\usetikzlibrary{arrows.meta,positioning,shapes.geometric,fit}
\usepackage[hidelinks]{hyperref}
\usepackage{enumitem}
\usepackage{caption}
\usepackage{float}
\setlist{nosep}
\definecolor{codebg}{RGB}{247,248,250}
\definecolor{codeblue}{RGB}{30,70,140}
\definecolor{codegreen}{RGB}{25,115,70}
\lstdefinestyle{python}{
  language=Python,
  basicstyle=\ttfamily\small,
  backgroundcolor=\color{codebg},
  keywordstyle=\color{codeblue}\bfseries,
  commentstyle=\color{codegreen},
  stringstyle=\color{red!65!black},
  numbers=left, numberstyle=\tiny\color{gray},
  breaklines=true, frame=single, rulecolor=\color{black!15},
  showstringspaces=false, columns=fullflexible
}
\title{Report 2: Detailed Walkthrough of the PyTorch Grad-CAM Implementation}
\author{Technical Study Notes}
\date{\today}
\begin{document}
\maketitle
\tableofcontents
\newpage
\section{Objective and complete code context}
This report explains the implementation line by line. The target is the final layer of the audio tower, and the scalar objective is the difference between the PD and HC next-token logits. The resulting explanation is a temporal Grad-CAM over audio-token positions.

\section{Select the target module}
\begin{lstlisting}[style=python]
target_layer = model.model.audio_tower.layers[-1]
\end{lstlisting}
\texttt{model} is the complete multimodal model. \texttt{model.model.audio\_tower} accesses its audio encoder. \texttt{layers[-1]} selects the last element of the audio tower's layer sequence. This is the last layer, not the first output of that layer.

The final layer is often selected because its representations are relatively high-level, but it is not universally the best layer. Different target layers can produce explanations at different abstraction levels and temporal resolutions.

\section{Create storage}
\begin{lstlisting}[style=python]
activations = []
gradients = []
\end{lstlisting}
The activation list stores the target layer's forward output. The gradient list is declared but is not used in the supplied implementation: the code obtains the gradient from \texttt{activation.grad}. It can therefore be removed unless the implementation is later changed to append gradients through a hook.

\section{Define the forward hook}
\begin{lstlisting}[style=python]
def forward_hook(module, inp, output):
    activation = output[0] if isinstance(output, tuple) else output
    activations.append(activation)
    activation.retain_grad()
\end{lstlisting}
PyTorch calls this function after the target module executes. Its arguments are:
\begin{itemize}
\item \texttt{module}: the module that has just run;
\item \texttt{inp}: the inputs passed to that module;
\item \texttt{output}: the module's returned value.
\end{itemize}

The conditional expression handles two possible output structures. If \texttt{output} is a tuple, \texttt{output[0]} is selected. Otherwise, the output itself is selected. This is an architectural assumption: inspect the output before relying on it. A tuple's first element is not inherently guaranteed to be a hidden-state tensor.

The selected tensor is appended to \texttt{activations}. The call \texttt{retain\_grad()} does not calculate a gradient immediately. It requests that PyTorch retain this intermediate tensor's gradient when a later backward pass computes it.

\section{Register the hook}
\begin{lstlisting}[style=python]
handle = target_layer.register_forward_hook(forward_hook)
\end{lstlisting}
This attaches the callback to the target module. Whenever the module executes in a forward pass, the callback runs. The returned handle is used to remove the hook after the explanation so that later forwards do not continue appending activations.

\section{Set evaluation mode and clear parameter gradients}
\begin{lstlisting}[style=python]
model.eval()
model.zero_grad(set_to_none=True)
\end{lstlisting}
\texttt{model.eval()} switches modules such as dropout and batch normalization into evaluation behavior. It does not disable autograd. This is important: Grad-CAM still requires gradients.

\texttt{model.zero\_grad(set\_to\_none=True)} clears previously accumulated parameter gradients. PyTorch accumulates gradients by default, so clearing them avoids mixing gradients from earlier operations. This call is distinct from retaining the activation gradient.

\section{Forward pass}
\begin{lstlisting}[style=python]
outputs = model(
    **inputs,
    use_cache=False,
    return_dict=True
)
\end{lstlisting}
The dictionary \texttt{inputs} is unpacked into keyword arguments. It may contain the Mel-spectrogram features and any other inputs required by the architecture. \texttt{use\_cache=False} disables autoregressive key-value caching, which is not needed for this gradient explanation. \texttt{return\_dict=True} requests a structured output with named fields such as \texttt{logits}.

During this forward pass, the target layer executes and the registered hook captures its activation. Because gradient tracking has not been disabled, PyTorch records the operations needed for backpropagation.

\section{Define the PD-versus-HC score}
\begin{lstlisting}[style=python]
score = (
    outputs.logits[0, -1, pd_ids[0]]
    - outputs.logits[0, -1, hc_ids[0]]
)
\end{lstlisting}
For logits shaped $(B,L,V)$, the indices select:
\begin{itemize}
\item batch index $0$: the first example;
\item sequence index $-1$: the final sequence position;
\item \texttt{pd\_ids[0]}: the selected PD token ID;
\item \texttt{hc\_ids[0]}: the selected HC token ID.
\end{itemize}
Thus,
\[
s=z_{\mathrm{PD}}-z_{\mathrm{HC}}.
\]
A positive score means the PD token logit is larger; a negative score means the HC token logit is larger. The score is a scalar, which is convenient for backpropagation.

The code assumes that each class is represented by the selected token ID. If labels tokenize into multiple tokens, a sequence-level scoring method may be required. Also verify that the final position is indeed the position whose next-token distribution represents the classification output.

\section{Backpropagate}
\begin{lstlisting}[style=python]
score.backward()
\end{lstlisting}
This computes derivatives of the scalar score with respect to tensors in the computation graph, including the retained target activation if it is connected to the score:
\[
G=\frac{\partial s}{\partial A}.
\]
The forward hook did not compute this derivative. It captured the activation and requested gradient retention; \texttt{backward()} computes the gradient afterward.

\section{Retrieve activation and gradient}
\begin{lstlisting}[style=python]
activation = activations[0]
grad = activation.grad

if grad is None:
    raise RuntimeError(
        "No gradient captured for the audio encoder layer."
    )
\end{lstlisting}
\texttt{activations[0]} retrieves the first captured activation. This is appropriate when the target module executes once during the forward pass. If the module is reused or called multiple times, inspect the list and select the intended invocation.

\texttt{activation.grad} retrieves the retained gradient. It should have the same shape as the activation. The explicit check catches cases where no gradient was retained. Possible causes include disabled gradient tracking, a broken or detached computation path, failure to call \texttt{retain\_grad()} on the relevant tensor, or a score that does not depend on the selected layer output.

\section{Compute channel weights}
\begin{lstlisting}[style=python]
weights = grad.mean(dim=1, keepdim=True)
\end{lstlisting}
Assume the activation and gradient shapes are $(1,T,D)$. Dimension $1$ is the audio-token position dimension. The mean produces
\[
\alpha_k=\frac{1}{T}\sum_{t=1}^{T}G_{t,k}.
\]
With \texttt{keepdim=True}, the result has shape $(1,1,D)$ rather than $(1,D)$. This allows broadcasting against the activation tensor $(1,T,D)$.

\section{Construct the temporal CAM}
\begin{lstlisting}[style=python]
cam = torch.relu(
    (weights * activation).sum(dim=-1)
)[0]
\end{lstlisting}
First, broadcasting multiplies each feature channel at each position by its channel weight:
\[
Q_{t,k}=\alpha_k A_{t,k}.
\]
Next, \texttt{sum(dim=-1)} sums over the hidden dimension:
\[
C_t=\sum_{k=1}^{D}\alpha_k A_{t,k}.
\]
Then ReLU gives
\[
L_t=\max(0,C_t).
\]
Finally, \texttt{[0]} removes the batch dimension. The result has shape $(T,)$. ReLU retains positive evidence for the selected score under this Grad-CAM construction and removes negative values.

\section{Move the CAM to NumPy}
\begin{lstlisting}[style=python]
cam = cam.detach().float().cpu().numpy()
\end{lstlisting}
\texttt{detach()} removes the tensor from the autograd graph for subsequent visualization. \texttt{float()} converts it to float32, \texttt{cpu()} moves it to CPU memory, and \texttt{numpy()} converts it to a NumPy array.

\section{Inspect the input feature dimensions}
\begin{lstlisting}[style=python]
input_features = inputs["input_features"][0]
mel_bins, time_frames = input_features.shape
\end{lstlisting}
The first operation selects the first example. The code assumes the resulting feature matrix is arranged as $(M,T_{\mathrm{in}})$: Mel bins by time frames. It stores these dimensions for later resizing and visualization.

\section{Reshape for one-dimensional interpolation}
\begin{lstlisting}[style=python]
cam_tensor = torch.tensor(
    cam, dtype=torch.float32
).view(1, 1, -1)
\end{lstlisting}
A one-dimensional CAM of length $T$ is reshaped to $(1,1,T)$. PyTorch's one-dimensional interpolation expects a tensor organized as batch, channels, and length. The \texttt{-1} asks PyTorch to infer the final dimension.

\section{Resize to the spectrogram time length}
\begin{lstlisting}[style=python]
cam_resized = torch.nn.functional.interpolate(
    cam_tensor,
    size=time_frames,
    mode="linear",
    align_corners=False
)[0, 0].numpy()
\end{lstlisting}
The output length is set to the input spectrogram's number of time frames. Linear interpolation estimates intermediate values between the original CAM samples. \texttt{align\_corners=False} selects the standard non-corner-aligned coordinate convention. Indexing with \texttt{[0,0]} removes batch and channel dimensions, yielding a vector of length \texttt{time\_frames}.

This operation changes display resolution; it does not create additional model evidence. The mapping assumes that proportional alignment between token positions and spectrogram frames is a reasonable approximation. Encoder strides, padding, chunking, or other transformations can make exact alignment more complicated.

\section{Broadcast the temporal map across Mel bins}
\begin{lstlisting}[style=python]
heatmap = np.repeat(
    cam_resized[np.newaxis, :],
    mel_bins,
    axis=0
)
\end{lstlisting}
\texttt{np.newaxis} changes a vector of shape $(T_{\mathrm{in}})$ into a row of shape $(1,T_{\mathrm{in}})$. Repeating that row $M$ times produces a heatmap of shape $(M,T_{\mathrm{in}})$.

Every frequency bin at the same time frame receives the same score. Consequently, this is a temporal map displayed over a two-dimensional spectrogram, not a frequency-resolved attribution map. It cannot establish that one Mel-frequency region is more important than another at a given time.

\section{Normalize}
\begin{lstlisting}[style=python]
if heatmap.max() > 0:
    heatmap /= heatmap.max()
\end{lstlisting}
If the maximum is positive, all values are divided by that maximum:
\[
H_{\mathrm{norm}}=\frac{H}{\max(H)}.
\]
The maximum becomes one. This makes the map convenient for visualization but makes values relative to that example's maximum. A normalized value of one in one recording is not necessarily comparable in absolute magnitude to one in another recording.

\section{Convert the input features for plotting}
\begin{lstlisting}[style=python]
features_np = input_features.detach().float().cpu().numpy()
\end{lstlisting}
This detaches the input features, converts them to float32, moves them to CPU, and converts them to NumPy. The resulting array can be passed to plotting functions alongside the heatmap.

\section{Tensor-shape summary}
The following shapes are illustrative; actual dimensions depend on the model and input.
\begin{center}
\begin{tabularx}{\linewidth}{>{\ttfamily}l X}
\toprule
\textnormal{Expression} & \textnormal{Illustrative shape / meaning}\\
\midrule
\texttt{activation} & $(1,T,D)$, e.g. $(1,200,1280)$\\
\texttt{grad} & $(1,T,D)$, same shape as activation\\
\texttt{weights} & $(1,1,D)$ after averaging over tokens\\
\texttt{cam} & $(T,)$ after summing features, ReLU, and removing batch\\
\texttt{cam\_tensor} & $(1,1,T)$ for 1-D interpolation\\
\texttt{cam\_resized} & $(T_{\mathrm{in}},)$\\
\texttt{heatmap} & $(M,T_{\mathrm{in}})$\\
\texttt{features\_np} & $(M,T_{\mathrm{in}})$\\
\bottomrule
\end{tabularx}
\end{center}

\section{Recommended hook cleanup}
The handle should be removed after the computation, even if an exception occurs:
\begin{lstlisting}[style=python]
handle = target_layer.register_forward_hook(forward_hook)
try:
    # forward pass, score, backward pass, and CAM computation
    ...
finally:
    handle.remove()
\end{lstlisting}
This avoids retaining a hook across repeated explanations. For repeated calls, also clear the activation list before each forward pass so old captures are not accidentally reused.

\section{Implementation checklist}
\begin{enumerate}
\item Confirm that the target layer is the intended final audio-tower layer.
\item Inspect the layer output type and shape; confirm that tuple element zero is the hidden state.
\item Ensure autograd is enabled and the activation is connected to the selected score.
\item Confirm that the class token IDs and final sequence position represent the intended decision.
\item Confirm that the module executes the expected number of times.
\item Remove the hook in a \texttt{finally} block.
\item Treat the repeated-row heatmap as temporal-only, not frequency-specific.
\end{enumerate}

\section{Conclusion}
The implementation captures the final audio layer's activation during the forward pass, defines a scalar PD-minus-HC logit score, and uses backpropagation to obtain the score gradient with respect to that activation. Gradient averaging creates channel weights; the weighted activation sum and ReLU produce one score per audio-token position. The remaining operations resize this temporal vector and display it over the Mel-spectrogram grid.

\end{document}
